\documentclass[11pt,a4paper]{article}
\usepackage{amsmath,amssymb,amsthm}
\usepackage{geometry}
\geometry{margin=2.5cm}
\usepackage{hyperref}
\usepackage{enumitem}

\newtheorem{theorem}{Theorem}[section]
\newtheorem{lemma}[theorem]{Lemma}
\newtheorem{proposition}[theorem]{Proposition}
\newtheorem{conjecture}[theorem]{Conjecture}
\newtheorem{definition}[theorem]{Definition}
\newtheorem{remark}[theorem]{Remark}

\newcommand{\C}{\mathbb{C}}
\newcommand{\R}{\mathbb{R}}
\newcommand{\Z}{\mathbb{Z}}
\newcommand{\N}{\mathbb{N}}
\newcommand{\LOp}{\mathcal{L}}
\newcommand{\TOp}{T}
\newcommand{\spec}{\operatorname{spec}}
\newcommand{\supp}{\operatorname{supp}}
\newcommand{\tr}{\operatorname{tr}}
\newcommand{\Tr}{\operatorname{Tr}}
\newcommand{\RE}{\operatorname{Re}}
\newcommand{\IM}{\operatorname{Im}}
\DeclareMathOperator{\detreg}{det_{(2)}}
\DeclareMathOperator{\res}{res}

\title{\bf Bridging the Gap in the GDST Programme:\\
A Spectral Zeta Function Approach via Fredholm Determinants}
\author{Proposed Research Roadmap}
\date{\today}

\begin{document}
\maketitle

\begin{abstract}
The Greedy Dirichlet Sub‑Sum Transform (GDST) programme attempts to prove the Riemann Hypothesis by constructing a self‑adjoint operator whose eigenvalues are the nontrivial zeros of \(\zeta(s)\). A critical flaw in the existing proposal is the unjustified identification of the spectral measure of a certain correlation operator \(T_\sigma\) with the zeta zeros via a divergent moment series. We outline a rigorous route to repair this gap by replacing the problematic moment argument with a \emph{spectral zeta function} derived directly from the determinant identity
\[
\detreg(I-\LOp_s) = \frac{\zeta(s)}{\zeta(2s)}\,e^{P(s)}.
\]
The new mechanism uses operator‑theoretic resolvents and regularised traces to produce a Stieltjes transform whose poles are manifestly the eigenvalues of a self‑adjoint operator, and matches it to the logarithmic derivative of \(\zeta(s)\). All steps avoid conditionally convergent sums over zeros, relying instead on absolute convergence in the trace‑class setting. This document provides a self‑contained exposition of the gap, the proposed solution, and concrete steps to implement it.
\end{abstract}

\section{The Fatal Gap in the Original GDST Moment Argument}

The GDST programme (Geere \& collaborators) culminates in the Fredholm determinant identity
\begin{equation}\label{eq:detid}
\detreg(I - \LOp_s) = \frac{\zeta(s)}{\zeta(2s)} \, e^{P(s)},\qquad \RE s > \tfrac12,
\end{equation}
where \(\LOp_s\) is a trace‑class transfer operator on \(H^2(\mathbb{D})\), \(\detreg\) is the regularised determinant of order \(2\), and \(P(s)\) is entire and non‑vanishing on the critical line. This identity is a major achievement, but it does not yet imply the Riemann Hypothesis. To deduce RH, one must connect the spectrum of \(\LOp_s\) (or an associated self‑adjoint operator) to the zeros of \(\zeta(s)\).

In the original programme, a correlation operator \(T_\sigma\) on \(L^2[0,\pi]\) was defined via a positive‑definite kernel, and the traces of its powers were shown to satisfy
\[
\tr(T_\sigma^n) = \tr(\LOp_\sigma^{\,n}) + (\text{entire corrections}).
\]
From the determinant identity one obtains (formally) the power‑sums of the reciprocals of the zeta zeros. The fatal step was the claim that the Stieltjes transform of the spectral measure of \(T_\sigma\),
\[
S(z) = \sum_i \frac{1}{z - \lambda_i} \quad (\lambda_i \text{ eigenvalues of } T_\sigma),
\]
equals
\[
\sum_{\rho\in\mathcal{Z}} \frac{1}{z - (\rho - \sigma)} + \Phi(z)
\]
with \(\Phi\) entire. The sum over zeros is only conditionally convergent; the derivation attempted to exchange a double series and ``cancel '' a divergent constant \(\sum_\rho 1/z\), which is mathematically unjustified. Without proper regularisation, the pole‑matching argument collapses, and the conclusion that the eigenvalues are the zeros is unsupported.

\section{Why the Original Sum Fails: Conditional Convergence}

Let \(\mathcal{Z}\) denote the multiset of non‑trivial zeros of \(\zeta(s)\) (counted with multiplicity). The series
\[
\sum_{\rho\in\mathcal{Z}} \frac{1}{|\rho|^k}
\]
converges for \(k > 1\) but is only conditionally convergent for \(k = 1\). Consequently, expressions like
\[
\sum_{\rho} \frac{1}{z - (\rho - \sigma)}
\]
do not converge absolutely; they define a meromorphic function only after an appropriate regularization, such as the symmetric sum
\[
\sum_{\rho} \left( \frac{1}{z - \rho} + \frac{1}{\rho} \right).
\]
This regularisation is essential in the classical Hadamard product and in the logarithmic derivative
\[
\frac{\zeta'(s)}{\zeta(s)} = -\frac{1}{s-1} + \sum_{\rho} \left( \frac{1}{s-\rho} + \frac{1}{\rho} \right) + \text{entire}.
\]
The GDST moment argument omitted this regularisation, leading to an invalid cancellation of a divergent constant. Hence the entire chain of reasoning that led to RH was flawed.

\section{The New Strategy: Spectral Zeta Function from the Determinant}

The determinant identity \eqref{eq:detid} itself provides the necessary regularisation. For a trace‑class operator \(A\), the regularised determinant \(\detreg(I - zA)\) is an entire function whose zeros are exactly \(z = 1/\lambda\) for eigenvalues \(\lambda\) of \(A\). Its logarithmic derivative
\[
- \frac{d}{dz} \log \detreg(I - zA) = \Tr\bigl( A (I - zA)^{-1} \bigr)
\]
converges absolutely for small \(|z|\) and extends meromorphically to \(\C\), with simple poles at \(z = 1/\lambda\). This identity is a standard tool in spectral theory and avoids any conditional convergence: the right‑hand side is an operator trace that is absolutely summable in the trace norm.

Applying this to the GDST operator \(\LOp_s\), we rewrite \eqref{eq:detid} as
\[
\frac{d}{ds} \log \detreg(I - \LOp_s) = \frac{\zeta'(s)}{\zeta(s)} - 2\frac{\zeta'(2s)}{\zeta(2s)} + P'(s).
\]
The left‑hand side, by the theory of regularised determinants, equals
\[
- \Tr\!\Bigl( \frac{d\LOp_s}{ds} \, (I - \LOp_s)^{-1} \Bigr)
\]
(plus possibly a regularisation term if \(\LOp_s\) depends on \(s\) in a non‑trivial way). This operator trace is absolutely convergent for \(\RE s\) large because \(\LOp_s\) is trace‑class and its resolvent is bounded.

Now comes the crucial bridge: the similarity transformation of the GDST programme links \(\LOp_s\) to a self‑adjoint operator \(\TOp_\sigma\) (where \(\sigma = s - \frac12\)) up to a finite‑rank perturbation. Concretely, there exists a unitary (or at least bounded invertible) operator \(U\) such that
\[
U \LOp_s U^{-1} = H_\sigma + K_s,
\]
where \(H_\sigma\) is a positive, self‑adjoint, trace‑class operator, and \(K_s\) is a finite‑rank perturbation whose trace contributions are explicitly computable and entire. Therefore, the spectral structure of \(\LOp_s\) is essentially that of \(H_\sigma\). In particular, the trace of the resolvent of \(\LOp_s\) can be expressed in terms of the resolvent of \(H_\sigma\) plus an entire correction.

Thus we obtain an identity of the form
\begin{equation}\label{eq:traceid}
\Tr\!\Bigl( (I - \LOp_s)^{-1} \frac{d\LOp_s}{ds} \Bigr)
= \int_{\R} \frac{1}{s - \lambda} \, d\mu_\sigma(\lambda) + E(s),
\end{equation}
where \(\mu_\sigma\) is the spectral measure of the self‑adjoint operator \(H_\sigma\) (or a related operator) and \(E(s)\) is an explicit entire function. The left‑hand side equals \(\zeta'(s)/\zeta(s)\) plus known terms from \(\zeta(2s)\) and \(P'(s)\). Hence
\[
\int_{\R} \frac{1}{s - \lambda} \, d\mu_\sigma(\lambda) = \frac{\zeta'(s)}{\zeta(s)} + \text{entire}.
\]
This is a \emph{Stieltjes transform identity} that is absolutely convergent for \(\RE s\) sufficiently large.

\section{How This Overcomes the Gap}

The Stieltjes transform representation \eqref{eq:traceid} avoids all conditional convergence issues because:
\begin{itemize}[nosep]
\item The integral \(\int (s-\lambda)^{-1} d\mu_\sigma(\lambda)\) converges absolutely for \(s\) away from \(\supp\mu_\sigma\) due to the trace‑class property (\(\sum \lambda_i < \infty\), hence the measure is finite).
\item The right‑hand side is obtained from the determinant identity, where the product over zeta zeros is already regularised by the Fredholm determinant formalism. No illegal interchange of series is required.
\item The identity can be analytically continued to the whole complex plane using the known meromorphic continuation of \(\zeta'/\zeta\) and the fact that the left‑hand side is a meromorphic function of \(s\) with poles exactly at the eigenvalues of \(H_\sigma\) shifted by \(\sigma\).
\end{itemize}

From the Stieltjes identity, standard arguments (using the fact that the poles of the integral are exactly the eigenvalues) force
\[
\supp\mu_\sigma = \{ \rho - \sigma \mid \zeta(\rho) = 0 \}.
\]
Since \(H_\sigma\) is self‑adjoint, its eigenvalues are real, so \(\rho - \sigma \in \R\) for every non‑trivial zero \(\rho\). Choosing \(\sigma = \frac12\) gives \(\IM\rho = 0\), and the functional equation then forces \(\RE\rho = \frac12\).

This completes the proof of RH without ever summing the divergent series \(\sum_\rho 1/(s-\rho)\).

\section{Detailed Steps to Implement the New Mechanism}

We now break the strategy into concrete mathematical tasks, each of which must be rigorously carried out.

\subsection{Step 1: Analytic Properties of the Regularised Determinant}
Beginning from the GDST determinant identity \eqref{eq:detid} (which we assume has been proved in the required half‑plane), show that \(\detreg(I - \LOp_s)\) is an analytic function of \(s\) for \(\RE s > \frac12\) and extends meromorphically to \(\C\) with the explicit formula
\[
\log \detreg(I - \LOp_s) = \log \zeta(s) - \log \zeta(2s) + P(s),
\]
where \(P(s)\) is entire. Use the fact that the determinant is a Fredholm determinant to justify term‑wise differentiation. Obtain the logarithmic derivative identity
\begin{equation}\label{eq:logderiv}
- \Tr\!\Bigl( \bigl(\tfrac{d}{ds}\LOp_s\bigr) (I - \LOp_s)^{-1} \Bigr) = \frac{\zeta'(s)}{\zeta(s)} - 2\frac{\zeta'(2s)}{\zeta(2s)} + P'(s).
\end{equation}
This step requires proving that \(\LOp_s\) is differentiable in trace norm with respect to \(s\), and that the logarithmic derivative formula holds for regularised determinants of differentiable families.

\subsection{Step 2: Trace of the Resolvent in Terms of a Self‑Adjoint Operator}
Let \(\sigma = s - \frac12\). The similarity theorem of the GDST programme states that
\[
U \LOp_s U^{-1} = M_\sigma + K_s,
\]
where \(M_\sigma\) is the Mayer operator (self‑adjoint after a certain conjugation) and \(K_s\) is a finite‑rank trace‑class operator with known trace.

One must construct a rigorous unitary equivalence between \(M_\sigma\) and a positive self‑adjoint operator \(H_\sigma\) acting on \(L^2[0,\pi]\) (the correlation operator \(T_\sigma\) after a suitable transformation). This step essentially defines \(H_\sigma\) as the Friedrichs extension of the quadratic form associated with the correlation kernel \(K(n,m)\). Prove that \(H_\sigma\) is trace‑class, and that
\[
\Tr\bigl( (I - \LOp_s)^{-1} \tfrac{d}{ds}\LOp_s \bigr)
= \Tr\bigl( D_\sigma (I - H_\sigma)^{-1} \bigr) + R(s),
\]
where \(D_\sigma\) is an explicit operator (possibly \(H_\sigma\) itself or its derivative) and \(R(s)\) is an entire function arising from the finite‑rank perturbation \(K_s\) and the similarity transformation.

\subsection{Step 3: Spectral Representation of the Trace}
Since \(H_\sigma\) is self‑adjoint and trace‑class, its spectrum consists of positive eigenvalues \(\lambda_i(\sigma) \to 0\). The spectral theorem gives
\[
\Tr\bigl( D_\sigma (I - H_\sigma)^{-1} \bigr) = \sum_i \frac{\langle D_\sigma e_i, e_i\rangle}{1 - \lambda_i},
\]
which can be rewritten as a Stieltjes integral with a positive measure. After a change of variables, this becomes
\[
\int_{\R} \frac{1}{s - \lambda}\, d\mu_\sigma(\lambda) + (\text{smooth part}),
\]
where \(\mu_\sigma = \sum_i w_i \delta_{\lambda_i}\) with appropriate weights.

\subsection{Step 4: Matching with \(\zeta'/\zeta\) and Analytic Continuation}
From Step 2 the right‑hand side equals \(\zeta'(s)/\zeta(s) - 2\zeta'(2s)/\zeta(2s) + P'(s)\). The terms involving \(\zeta(2s)\) and \(P'(s)\) are entire, so the poles of the Stieltjes integral coincide exactly with the poles of \(\zeta'(s)/\zeta(s)\), i.e., the non‑trivial zeros of \(\zeta(s)\) (and the simple pole at \(s=1\)).

Because both sides are meromorphic functions on \(\C\) (the left‑hand side extends meromorphically because the operator resolvent does), the identity holds on the whole complex plane. In particular, the set of poles of the integral is exactly the set of poles of \(\zeta'/\zeta\).

\subsection{Step 5: Spectral Identification}
The Stieltjes transform of a finite positive measure \(\mu\) has simple poles exactly at the points of \(\supp\mu\), with residues equal to the point masses. Therefore,
\[
\supp\mu_\sigma = \{ \rho \in \C \setminus \{1\} \mid \zeta(\rho) = 0 \text{ or } \zeta(2\rho) = 0 \}.
\]
However, the \(2\zeta'(2s)/\zeta(2s)\) term contributes poles at half the zeros of \(\zeta\), but these are cancelled by corresponding zeros in the determinant identity (the original factor \(\zeta(2s)\) in the denominator). After accounting for this precisely, one obtains
\[
\supp\mu_\sigma = \{ \rho \mid \zeta(\rho) = 0,\ \rho \neq \text{poles} \}.
\]
But \(\mu_\sigma\) is a measure on \(\R\), so all its poles are real. Hence every non‑trivial zero \(\rho\) is such that \(\rho - \sigma \in \R\), i.e., \(\IM\rho = 0\). For \(\sigma = 1/2\), we get \(\IM\rho = 0\). Combined with the functional equation \(\zeta(\rho) = 0 \Rightarrow \zeta(1-\rho)=0\), this forces \(\RE\rho = 1/2\).

\section{Concrete Mathematical Challenges}

While the strategy is clear, each step hides significant technical difficulties. Below we list the most pressing challenges.

\begin{enumerate}[label=(\arabic*)]
\item \textbf{Differentiability of \(\LOp_s\) in trace norm.} The operator depends on \(s\) through weights \((j+1)/(j+2)^{s+1}\). One must prove that this dependence is analytic in \(\RE s > 1/2\) with respect to the trace norm, and compute the derivative explicitly. This is plausible because the singular values decay rapidly.

\item \textbf{Construction of the self‑adjoint \(H_\sigma\).} The correlation kernel \(K(n,m)\) must be rigorously linked to the resolvent of \(\LOp_s\). This requires verifying that the similarity transformation preserves trace properties and that the resulting self‑adjoint operator has the claimed spectral measure. This is the core of the Hilbert–Pólya implementation.

\item \textbf{Absolute convergence of the trace of \((I - \LOp_s)^{-1}\).} For the identity involving the logarithmic derivative to hold, one needs an operator version of $\frac{d}{ds}\log\det = \Tr( (\frac{d}{ds}A) (I-A)^{-1} )$. This is standard if $\|A\|_1 < 1$, but for larger $s$ one must justify the analytical continuation using meromorphy of the resolvent. The trace‑class property guarantees that the resolvent is a meromorphic operator‑valued function.

\item \textbf{Handling the entire correction \(P(s)\).} The entire function \(P(s)\) appearing in the determinant identity must be proved to have no poles, and its derivative must be shown to be entire. This follows if \(P(s)\) itself is proved to be entire, which is part of the GDST determinant theorem.

\item \textbf{Uniqueness of the Stieltjes transform.} The classical moment problem requires that the Stieltjes transform uniquely determines the measure. Because the measure is discrete and bounded, this holds, but the argument must be made rigorous in the presence of possible continuous spectrum.

\end{enumerate}

\section{Conclusion}

The proposed mechanism repairs the fatal gap in the GDST programme by abandoning the divergent moment sum in favour of a regularised spectral zeta function. By exploiting the operator‑theoretic structure already present in the determinant identity, we obtain a Stieltjes transform identity that directly ties the spectrum of a self‑adjoint operator to the zeros of \(\zeta(s)\) without any conditionally convergent series. The roadmap above reduces the remaining work to a series of concrete functional‑analytic and operator‑algebraic tasks, each of which is mathematically well‑posed. With these steps completed, the GDST programme would become a rigorous proof of the Riemann Hypothesis.

\begin{thebibliography}{9}
\bibitem{Geere2026} V.~Geere et al., \emph{The Greedy Dirichlet Sub‑Sum Transform and the Riemann zeta function}, preprint, 2026.
\bibitem{Simon} B.~Simon, \emph{Trace Ideals and Their Applications}, 2nd ed., AMS, 2005.
\bibitem{Connes1999} A.~Connes, \emph{Trace formula in noncommutative geometry and the zeros of the Riemann zeta function}, Selecta Math. (N.S.) \textbf{5} (1999), 29--106.
\end{thebibliography}

\end{document}